% Options for packages loaded elsewhere
\PassOptionsToPackage{unicode}{hyperref}
\PassOptionsToPackage{hyphens}{url}
%
\documentclass[
]{article}
\usepackage{amsmath,amssymb}
\usepackage{iftex}
\ifPDFTeX
  \usepackage[T1]{fontenc}
  \usepackage[utf8]{inputenc}
  \usepackage{textcomp} % provide euro and other symbols
\else % if luatex or xetex
  \usepackage{unicode-math} % this also loads fontspec
  \defaultfontfeatures{Scale=MatchLowercase}
  \defaultfontfeatures[\rmfamily]{Ligatures=TeX,Scale=1}
\fi
\usepackage{lmodern}
\ifPDFTeX\else
  % xetex/luatex font selection
\fi
% Use upquote if available, for straight quotes in verbatim environments
\IfFileExists{upquote.sty}{\usepackage{upquote}}{}
\IfFileExists{microtype.sty}{% use microtype if available
  \usepackage[]{microtype}
  \UseMicrotypeSet[protrusion]{basicmath} % disable protrusion for tt fonts
}{}
\makeatletter
\@ifundefined{KOMAClassName}{% if non-KOMA class
  \IfFileExists{parskip.sty}{%
    \usepackage{parskip}
  }{% else
    \setlength{\parindent}{0pt}
    \setlength{\parskip}{6pt plus 2pt minus 1pt}}
}{% if KOMA class
  \KOMAoptions{parskip=half}}
\makeatother
\usepackage{xcolor}
\usepackage[margin=1in]{geometry}
\usepackage{longtable,booktabs,array}
\usepackage{calc} % for calculating minipage widths
% Correct order of tables after \paragraph or \subparagraph
\usepackage{etoolbox}
\makeatletter
\patchcmd\longtable{\par}{\if@noskipsec\mbox{}\fi\par}{}{}
\makeatother
% Allow footnotes in longtable head/foot
\IfFileExists{footnotehyper.sty}{\usepackage{footnotehyper}}{\usepackage{footnote}}
\makesavenoteenv{longtable}
\usepackage{graphicx}
\makeatletter
\def\maxwidth{\ifdim\Gin@nat@width>\linewidth\linewidth\else\Gin@nat@width\fi}
\def\maxheight{\ifdim\Gin@nat@height>\textheight\textheight\else\Gin@nat@height\fi}
\makeatother
% Scale images if necessary, so that they will not overflow the page
% margins by default, and it is still possible to overwrite the defaults
% using explicit options in \includegraphics[width, height, ...]{}
\setkeys{Gin}{width=\maxwidth,height=\maxheight,keepaspectratio}
% Set default figure placement to htbp
\makeatletter
\def\fps@figure{htbp}
\makeatother
\setlength{\emergencystretch}{3em} % prevent overfull lines
\providecommand{\tightlist}{%
  \setlength{\itemsep}{0pt}\setlength{\parskip}{0pt}}
\setcounter{secnumdepth}{-\maxdimen} % remove section numbering
\ifLuaTeX
  \usepackage{selnolig}  % disable illegal ligatures
\fi
\usepackage{bookmark}
\IfFileExists{xurl.sty}{\usepackage{xurl}}{} % add URL line breaks if available
\urlstyle{same}
\hypersetup{
  hidelinks,
  pdfcreator={LaTeX via pandoc}}

\author{}
\date{\vspace{-2.5em}}

\begin{document}

\subsection{Stoppage Time --- counterfactual model
review}\label{stoppage-time-counterfactual-model-review}

\textbf{What this is.} This is a counterfactual estimate over 314
already-played matches: \emph{if the stoppage minutes that were truly
owed but never played had actually been played, in how many of those
matches would at least one additional goal have been score? Within this
subset, how many results would have flipped (win/draw/loss)}

This document is meant to be read start to finish with no outside
material. Standard probability and bootstrap machinery is assumed; every
football- or project-specific term is defined where it first appears.
There are 7 sections. \textbf{Sections 2 and 4 are the parts I most want
scrutinized} --- 2 is where the data's ambiguity is resolved by
calibration, and 4 is where I explain the assumptions.

\begin{center}\rule{0.5\linewidth}{0.5pt}\end{center}

\subsection{1. The claim}\label{the-claim}

\textbf{In plain English.} Hold each match exactly as it was actually
played. Then \emph{append} the stoppage minutes the referee should have
added but didn't, and ask whether at least one more goal would plausibly
have landed in those appended minutes. Averaged over all 314 matches,
that probability is \textbf{X\% = 23.6\%}:

\emph{If stoppage time were measured and awarded in full, 23.6\% of the
314 matches would have finished with a different scoreline} --- at least
one extra goal somewhere in the omitted added time. This is a mean of
per-match probabilities, computed in closed form.

\begin{longtable}[]{@{}
  >{\raggedright\arraybackslash}p{(\columnwidth - 2\tabcolsep) * \real{0.5000}}
  >{\raggedright\arraybackslash}p{(\columnwidth - 2\tabcolsep) * \real{0.5000}}@{}}
\toprule\noalign{}
\begin{minipage}[b]{\linewidth}\raggedright
Quantity
\end{minipage} & \begin{minipage}[b]{\linewidth}\raggedright
Value
\end{minipage} \\
\midrule\noalign{}
\endhead
\bottomrule\noalign{}
\endlastfoot
\textbf{Headline --- different scoreline (central)} & \textbf{23.6\%} \\
95\% bootstrap confidence interval (sampling) & \textbf{{[}20.6\%,
27.4\%{]}} \\
Specification band --- one modeling choice varied at a time &
\textbf{21.1\% -- 26.1\%} \\
Full envelope --- all modeling choices varied jointly & \textbf{18.6\%
-- 27.3\%} \\
\textbf{Different \emph{result}} --- winner/draw actually changes &
\textbf{12.1\% {[}10.6\%, 14.2\%{]}} \\
Second-half-only variant (comparison, not the headline) & 16.0\%
{[}14.0\%, 18.5\%{]} \\
\end{longtable}

\emph{All figures are for the full match (first- and second-half
stoppage combined), pooled across all 314 matches. The two bands are
different objects and are kept apart throughout: the bootstrap CI is
\textbf{sampling} uncertainty; the specification band/envelope is how
X\% moves as \textbf{modeling choices} change (§5).}

\textbf{``Different scoreline'' vs ``different result'' --- keep these
apart.} ``≥1 extra goal anywhere'' (a different \emph{scoreline},
23.6\%) is strictly weaker than ``the winner or the draw actually
changes'' (a different \emph{result}, 12.1\%). Of the 314 matches, 95
were decided by two goals or more at the 90-minute mark; we treat those
as unflippable, so they can change scoreline but not result. Conflating
the two would overstate the claim, so the two numbers are always
reported separately.

\begin{center}\rule{0.5\linewidth}{0.5pt}\end{center}

\subsection{2. The data, its gaps, and how we resolve
them}\label{the-data-its-gaps-and-how-we-resolve-them}

\emph{This is the section where I'm probably most curious for your
input. The model's single largest judgment call lives here: the source
data has irreducibly ambiguous gaps, and the way we resolve them is to
calibrate against one external dataset (Nate Silver, World Cup 2018) and
then extrapolate to five other tournaments. The question I am really
asking is whether that calibrate-then-extrapolate move is sound.}

\textbf{The matches.} 314 matches from six major international
tournaments:

\begin{longtable}[]{@{}lll@{}}
\toprule\noalign{}
Tournament & Era & Matches \\
\midrule\noalign{}
\endhead
\bottomrule\noalign{}
\endlastfoot
FIFA World Cup 2018 & PRE & 64 \\
UEFA Euro 2020 & PRE & 51 \\
FIFA World Cup 2022 & POST & 64 \\
UEFA Euro 2024 & POST & 51 \\
Copa América 2024 & POST & 32 \\
Africa Cup of Nations 2023 & POST & 52 \\
\textbf{Total} & & \textbf{314} (PRE 115, POST 199) \\
\end{longtable}

A note on PRE vs POST: A 2022 IFAB directive (detail in §5) splits these
into a ``before'' (PRE, 115) and ``after'' (POST, 199) cohort by
\emph{how much added time officials were instructed to show}. Basically,
the instructions for the refs changed, and as a result, they added more
stoppage time. Interestingly, goals per live-minute remained
statistically indistinguishable across the split (§5), so the model
\textbf{pools all six tournaments} for its scoring rates. The PRE/POST
split matters only as a robustness axis --- it is not the headline, and
not what I need reviewed.

\textbf{Where the numbers come from: StatsBomb event data.} Every figure
traces to StatsBomb's open event data, which is \textbf{hand-coded}: an
analyst watches the match video and logs every on-ball event --- pass,
shot, tackle, throw-in, foul, card, substitution --- timestamped to the
second. During active play that is an event every few seconds. When the
ball is dead --- injury, goal celebration, substitution, VAR review ---
nothing is logged, so a \textbf{gap} opens in the timestamp stream. We
use only events, scores, and timestamps: no expected-goals (xG) model,
no player-tracking (``360'') data.

Because the data is event-driven, dead time is visible \emph{only} as
one of these gaps --- and the gaps are not all equal:

\begin{itemize}
\tightlist
\item
  \textbf{Unambiguous gaps} carry an explicit out-of-play marker on the
  last logged touch --- the ball is flagged \emph{out}, or a throw-in /
  goal kick / corner / free kick follows. We know the clock should have
  stopped, and roughly for how long.
\item
  \textbf{Data Gaps} are stretches of ≥ 20 seconds with no events
  \emph{and no out-of-play marker}. From the data alone we cannot tell
  whether the ball was genuinely dead (a stoppage we \emph{should}
  count) or simply being moved slowly --- often off-camera --- in open
  play (live time we should \emph{not} count).
\end{itemize}

\textbf{A real Data Gap --- the whole problem in one example.}
Austria--France, Euro 2024, second half. These are two consecutive
logged touches, with nothing recorded in between:

\begin{longtable}[]{@{}
  >{\raggedright\arraybackslash}p{(\columnwidth - 8\tabcolsep) * \real{0.2000}}
  >{\raggedright\arraybackslash}p{(\columnwidth - 8\tabcolsep) * \real{0.2000}}
  >{\raggedright\arraybackslash}p{(\columnwidth - 8\tabcolsep) * \real{0.2000}}
  >{\raggedright\arraybackslash}p{(\columnwidth - 8\tabcolsep) * \real{0.2000}}
  >{\raggedright\arraybackslash}p{(\columnwidth - 8\tabcolsep) * \real{0.2000}}@{}}
\toprule\noalign{}
\begin{minipage}[b]{\linewidth}\raggedright
Match clock
\end{minipage} & \begin{minipage}[b]{\linewidth}\raggedright
Event type
\end{minipage} & \begin{minipage}[b]{\linewidth}\raggedright
Play pattern
\end{minipage} & \begin{minipage}[b]{\linewidth}\raggedright
Ball flagged out?
\end{minipage} & \begin{minipage}[b]{\linewidth}\raggedright
Off-camera?
\end{minipage} \\
\midrule\noalign{}
\endhead
\bottomrule\noalign{}
\endlastfoot
50:12 & Ball Recovery (France) & Regular Play & no & no \\
50:38 & Pass (France) & Regular Play & no & yes \\
\end{longtable}

\emph{Label definitions, for a reader without the schema: \textbf{Ball
Recovery / Pass} are StatsBomb on-ball event types; \textbf{Regular
Play} means the touch happened in open play, not from a restart;
\textbf{Ball flagged out} is StatsBomb's explicit marker that the ball
had left the field of play; \textbf{Off-camera} flags a touch the
analyst coded while the broadcast had cut away from the field.}

Twenty-six seconds pass between these two touches. France has the ball
at 50:12 in open play, the ball is \textbf{not} flagged out, and the
next logged touch --- same possession, 26 seconds later --- is coded
off-camera. Was the ball dead for those 26 seconds (an untracked knock
or delay we \emph{should} credit as owed stoppage), or alive the whole
time (France simply circulating the ball while the camera was
elsewhere)? \textbf{Nothing in the data answers this.} Gaps like it
total roughly \textbf{8.4 minutes per match} --- far too much to credit
or discard wholesale.

\textbf{How we resolve the Data Gaps.} The method is the following:

\begin{enumerate}
\def\labelenumi{\arabic{enumi}.}
\tightlist
\item
  \textbf{Credit only marker-confirmed gaps.} A Data Gap is credited as
  dead \emph{only} when its leading edge carries an out-of-play marker
  --- i.e.~the data independently confirms the ball had already left
  play. The \textasciitilde three-quarters of Data Gaps that are
  \emph{unmarked} (like the example above) are \textbf{dropped}.
  Crediting them all instead produces gross over-counts: the worst
  single match, Germany--Sweden 2018, balloons to 19.7 owed minutes and
  agreement with external ground truth collapses.
\item
  \textbf{Add back one small residual constant.} Dropping every unmarked
  gap \emph{under}-counts. So we add a single flat per-match residual
  (24.2 s, identical for every match), fit once so the estimator's
  \emph{mean} lands on the external ground truth for 2018 (next
  paragraph). It is a constant, not a per-gap guess.
\end{enumerate}

\textbf{Does it work? Calibration against Nate Silver's hand-measured
2018 data.} Our one independent ground truth is Nate Silver's stopwatch
measurement of all 32 World Cup 2018 matches (then at FiveThirtyEight),
reporting per match the stoppage that \emph{should} have been added
(``expected'') and the stoppage that \emph{was} played (``actual'').
Building the owed-stoppage estimator one layer at a time and scoring
each against his ``expected'' column:

\begin{longtable}[]{@{}
  >{\raggedright\arraybackslash}p{(\columnwidth - 6\tabcolsep) * \real{0.2500}}
  >{\raggedright\arraybackslash}p{(\columnwidth - 6\tabcolsep) * \real{0.2500}}
  >{\raggedright\arraybackslash}p{(\columnwidth - 6\tabcolsep) * \real{0.2500}}
  >{\raggedright\arraybackslash}p{(\columnwidth - 6\tabcolsep) * \real{0.2500}}@{}}
\toprule\noalign{}
\begin{minipage}[b]{\linewidth}\raggedright
Owed-stoppage estimator
\end{minipage} & \begin{minipage}[b]{\linewidth}\raggedright
Correlation r
\end{minipage} & \begin{minipage}[b]{\linewidth}\raggedright
MAE (min)
\end{minipage} & \begin{minipage}[b]{\linewidth}\raggedright
mean (min)
\end{minipage} \\
\midrule\noalign{}
\endhead
\bottomrule\noalign{}
\endlastfoot
identifiable dead time only (celebration / sub / card / injury) & 0.655
& 5.72 & 7.53 \\
+ excess over a standard restart allowance & 0.754 & 4.28 & 9.03 \\
+ marker-confirmed Data Gaps & 0.825 & 2.49 & 12.76 \\
+ residual constant (final estimator) & \textbf{0.825} & \textbf{2.44} &
12.60 \\
\end{longtable}

The final estimator tracks Nate's hand-measured ``expected'' at
\textbf{r = 0.825} (MAE 2.44 min) and averages \textbf{12.60 min/match
against his 13.16} --- within half a minute. The cross-check on the
\emph{other} column is what makes this credible: our independently
measured stoppage-\emph{played} clock matches his ``actual'' almost
exactly (\textbf{r = 0.992, MAE 0.13 min}). So the owed-minus-played gap
our pipeline produces reproduces the gap an outside observer found by
hand --- Nate's own numbers show \textbf{13.16 owed vs 6.98 played, a
6.2-minute (1.9×) shortfall} (Figure 1).

\emph{Figure 1. Validation against Nate Silver's independent,
hand-measured World Cup 2018 data, 32 matches. Left: our owed-stoppage
estimator vs his ``expected'' (should-have-been-added) minutes --- r =
0.825, MAE 2.44; the two labelled points are the model's hardest cases.
Right: our stoppage-played clock vs his ``actual'' (actually-played)
minutes --- r = 0.992, essentially exact. The dashed line is perfect
agreement; the near-exact right panel is what makes the left panel's gap
credible.}

\textbf{Two things to weigh --- this is the crux.} \emph{(i) Treating
the calibrated level as a point.} If the Data Gap credit is instead
swept from ``credit none'' to ``credit all,'' the headline X\% moves
from \textbf{10.8\% to 37.3\%} --- a wider swing than every other
modeling choice combined. We argue that ``none'' and ``all'' are
\emph{known-biased bounds}, not honest uncertainty, so we report the
calibrated middle as a single number and propagate only its calibration
error (the MAE) through the confidence interval (§5). That is a
defensible position, but it is the one decision the headline's tightness
rests on. \emph{(ii) Calibrate-then-extrapolate.} The entire calibration
above is fit on \textbf{World Cup 2018 only} --- the lone tournament
with independent ground truth --- and the fitted constants (restart
allowances, residual, marker rule) are then applied \textbf{unchanged}
to the other five tournaments, where estimated owed time runs
\textasciitilde1.5--2× higher (POST owed ≈ 19.4 min vs 2018's 12.6).
Whether a 2018-calibrated structure legitimately transfers to Copa
América, AFCON, and the rest is, in my judgment, the model's most
exposed assumption --- and the single thing I would most value your read
on.

\textbf{The perturbation, sized.} With the owed-stoppage estimator in
hand, for each match we have the stoppage \emph{owed}
(\texttt{T\_true}), the stoppage \emph{actually played}
(\texttt{T\_play}, whistle-to-whistle added time, read off the event
clock almost exactly), and their positive difference --- the
\emph{omitted} minutes \texttt{O} we inject counterfactually:

\begin{longtable}[]{@{}ll@{}}
\toprule\noalign{}
Per match (minutes) & All 314 \\
\midrule\noalign{}
\endhead
\bottomrule\noalign{}
\endlastfoot
Stoppage owed, \texttt{T\_true} & 17.3 \\
Stoppage actually played, \texttt{T\_play} & 8.9 \\
\textbf{Omitted (owed − played), \texttt{O}} & \textbf{8.6} \\
of which \emph{live} ball-in-play (after §4) & 5.1 \\
\end{longtable}

Omitted time is \textbf{positive in 97\% of matches} --- not a marginal
perturbation. By tournament, estimated owed time is Copa América 25.3,
AFCON 23.5, World Cup 2022 17.5, Euro 2020 16.3, Euro 2024 15.6, World
Cup 2018 12.6 min --- the \textasciitilde2× spread over the 2018 anchor
that drives the extrapolation worry above. One further thin spot to flag
for later: the first-half stoppage goal rate, a model input, rests on
just \textbf{23 goals} (§5).

\begin{center}\rule{0.5\linewidth}{0.5pt}\end{center}

\subsection{3. Estimand \&
identification}\label{estimand-identification}

\textbf{The counterfactual, formally.} For match
\texttt{i\ ∈\ \{1..314\}} and half-window \texttt{h\ ∈\ \{1H,\ 2H\}} (1H
= first half, 2H = second half):

\begin{verbatim}
T_true(i,h)   estimated true stoppage owed, in minutes (the estimator of §4.2)
T_play(i,h)   stoppage actually played = (period_end − 45:00 or 90:00), in minutes (measured)
O(i,h)        = max(0, T_true − T_play)        omitted (owed-but-unplayed) minutes
s(i,h)        = ball-in-play fraction of the window   ("live share"; measured, §4.1)
ℓ(i,h)        = O(i,h) · φ(s)                  omitted LIVE minutes
                  φ(s) = s                       (base model)
                  φ(s) = s · (1 + z(1−s))        (with the gross-up of §4.3)

λ_1H   = 0.0478   goals per live-minute, either team, in 1st-half stoppage   (n = 23 goals)
λ̄_2H(i)         = the 2nd-half rate after decay (§4.1): ramps from 0.0816 toward 0.0427

μ(i)   = λ_1H · ℓ(i,1H) + λ̄_2H(i) · ℓ(i,2H)   expected extra goals (either team)
P_change(i) = 1 − exp(−μ(i))                   probability of ≥1 extra goal (Poisson)
X%     = (1/314) · Σ_i P_change(i)
\end{verbatim}

The metric is \textbf{deterministic}; a random seed enters only the
confidence-interval bootstrap (§5). λ is the \textbf{two-team} rate
(goals by \emph{either} side per live-minute), so \texttt{μ\ =\ λ·ℓ} is
directly the expected \emph{count} of extra goals. (An earlier version
used a per-team rate and had to multiply by two to recover the match
rate; using the two-team rate from the start removes that factor-of-2
ambiguity.)

\textbf{Outcome-flip} replaces \texttt{P\_change(i)} with \texttt{Q(i)},
keyed on the score state at 90 minutes: level → \texttt{1\ −\ exp(−μ)}
(any goal breaks the tie); one-goal lead → \texttt{1\ −\ exp(−μ/2)}
(only the trailing team flips it, at half the two-team rate);
two-goal-or-more lead → 0. State counts at 90′: level 98, one-goal lead
121, two-or-more 95 --- so 219 of 314 are flip-eligible.

\textbf{What is observed vs imputed.}

\begin{longtable}[]{@{}
  >{\raggedright\arraybackslash}p{(\columnwidth - 4\tabcolsep) * \real{0.3333}}
  >{\raggedright\arraybackslash}p{(\columnwidth - 4\tabcolsep) * \real{0.3333}}
  >{\raggedright\arraybackslash}p{(\columnwidth - 4\tabcolsep) * \real{0.3333}}@{}}
\toprule\noalign{}
\begin{minipage}[b]{\linewidth}\raggedright
Component
\end{minipage} & \begin{minipage}[b]{\linewidth}\raggedright
Status
\end{minipage} & \begin{minipage}[b]{\linewidth}\raggedright
How obtained
\end{minipage} \\
\midrule\noalign{}
\endhead
\bottomrule\noalign{}
\endlastfoot
Goals, scoreline, score state at 45′/90′ & \textbf{Observed} & StatsBomb
events \\
\texttt{T\_play} stoppage actually played & \textbf{Observed}
(near-exact) & whistle-to-whistle from the event clock; r = 0.992 vs
Nate's ``actual'' (§6) \\
Ball-in-play / live share \texttt{s} & \textbf{Observed}, calibrated &
gap method (§4.1); pooled World Cup 2022 within 24s of Opta's published
figure \\
λ (stoppage \& open-play goal rates) & \textbf{Observed} & goals ÷
live-minutes, directly counted \\
\texttt{T\_true} stoppage owed & \textbf{Imputed} & the estimator of
§4.2, calibrated to Nate's ``expected'' (r = 0.825) \\
Extra goals in omitted minutes & \textbf{Imputed} & Poisson with λ̄
applied to ℓ \\
\end{longtable}

\textbf{Why there is no train/test split, and what stands in for
out-of-sample validity.} The estimand is a descriptive functional of a
\emph{fixed, fully-observed} population --- all 314 matches, not a
sample drawn from a super-population we forecast into. There is no
held-out outcome to predict and no future to backtest, so
cross-validation is undefined. Validity rests instead on (a) the
\textbf{inputs being unbiased} and (b) the \textbf{transfer assumptions
being defensible}. The only genuine out-of-sample anchoring is external,
and it is on the \emph{inputs}, not the output:

\begin{itemize}
\tightlist
\item
  the \textbf{owed-stoppage estimator} is calibrated to Nate Silver's
  hand-measured World Cup 2018 ``expected'' (should-be-added) column ---
  r = 0.825, mean absolute error 2.44 min, n = 32 matches;
\item
  the \textbf{ball-in-play / live share} is calibrated to Opta's
  published World Cup 2022 figure of 58:04 of ball-in-play per match (we
  reconstruct 57:40, 24s lower than the benchmark);
\item
  the \textbf{stoppage-played clock} matches Nate's ``actual''
  (actually-played) column almost exactly --- r = 0.992, MAE 0.13 min.
\end{itemize}

There is \textbf{no external check on X\% itself} --- see §6 and §7.

\textbf{Identifying assumptions, and where each breaks.}

\begin{longtable}[]{@{}
  >{\raggedright\arraybackslash}p{(\columnwidth - 4\tabcolsep) * \real{0.3333}}
  >{\raggedright\arraybackslash}p{(\columnwidth - 4\tabcolsep) * \real{0.3333}}
  >{\raggedright\arraybackslash}p{(\columnwidth - 4\tabcolsep) * \real{0.3333}}@{}}
\toprule\noalign{}
\begin{minipage}[b]{\linewidth}\raggedright
\#
\end{minipage} & \begin{minipage}[b]{\linewidth}\raggedright
Assumption
\end{minipage} & \begin{minipage}[b]{\linewidth}\raggedright
Breaks if\ldots{}
\end{minipage} \\
\midrule\noalign{}
\endhead
\bottomrule\noalign{}
\endlastfoot
A1 & \textbf{No re-simulation.} Appending omitted minutes does not
change the minutes actually played; each extra goal is an independent
increment toward ``≥1 extra goal.'' & \ldots game state matters --- a
first-half or early extra goal would change subsequent play. Internally
consistent for the \emph{scoreline} metric (we only need ≥1), but
``ended differently'' is not a full re-play. \\
A2 & \textbf{Productivity transfer.} Omitted minutes score at λ̄ --- the
observed in-stoppage rate, decayed toward the open-play floor. &
\ldots omitted minutes are systematically \emph{less} productive than
even open play (they are, after all, the minutes a referee declined to
play). The floor is the open-play rate, so even maximal decay never goes
below match-average pace. \\
A3 & \textbf{True-stoppage identification.} \texttt{T\_true} =
identifiable dead time + marker-gated Data Gaps + a small residual,
pinned by calibration to Nate's ``expected'' (§2). & \ldots POST behaves
unlike World Cup 2018. Calibration exists for World Cup 2018 only; POST
owed time (\textasciitilde19.4 min) is a \textasciitilde1.5×
extrapolation of the 2018 anchor (12.6), with all constants frozen on
2018. \\
A4 & \textbf{Rate exchangeability.} Goals-per-live-minute is the same
PRE/POST; omitted minutes draw the pooled in-stoppage rate. & \ldots the
rule change altered scoring pace \emph{per live minute} (it changed the
board, not the football). The PRE/POST sweep brackets this --- 22.6\% vs
26.1\%, within Poisson noise. \\
A5 & \textbf{Poisson within window.} Extra goals \textasciitilde{}
Poisson(μ); \texttt{P(≥1)\ =\ 1\ −\ e\^{}(−μ)}. & \ldots goals
over-disperse/cluster. But \texttt{P(≥1)} depends only on the
pre-first-goal hazard, so mild clustering is second-order; the observed
rate already embeds some second-goal inflation (8 of 73
second-half-stoppage goals are a 2nd+ goal in the same window). \\
A6 & \textbf{Live-share neutrality.} In the base model \texttt{s}
cancels in μ (it scales λ up and ℓ down equally), so μ ≈
goals-per-\emph{clock}-min × omitted-\emph{clock}-min. & \ldots only
through the gross-up --- the one place \texttt{s} enters non-trivially.
Off → geometric moves X\% by ≤3 pts (§4.3). \\
\end{longtable}

\begin{center}\rule{0.5\linewidth}{0.5pt}\end{center}

\subsection{4. The load-bearing
assumptions}\label{the-load-bearing-assumptions}

This is the part worth the most scrutiny, alongside §2. Three
assumptions carry the result below; a fourth --- how the ambiguous Data
Gaps are resolved --- is in §2. Each is stated with its rationale and
the conditions under which it breaks.

\subsubsection{4.1 Productivity transfer and the late-game
decay}\label{productivity-transfer-and-the-late-game-decay}

\textbf{Takeaway.} Every omitted minute must be assigned a goal rate.
Late-game minutes are observably more productive than average --- but we
do \emph{not} assume the newly-added minutes inherit the full late-game
premium. Their rate \textbf{decays} from the observed in-stoppage rate
back toward the open-play rate as the hypothetical added window grows.
The speed of that decay (its half-life) is the reported productivity
band.

First, the motivation. Goals do not arrive at a constant rate through a
match: the per-live-minute scoring rate climbs through the regulation
buckets, and the two stoppage windows are higher still (Figure 2).

\emph{Figure 2. Goal productivity (goals per live-minute, either team)
pooled across all six tournaments. Blue bars are regulation buckets; the
two red bars are the rates inside first- and second-half stoppage ---
the windows the model actually prices. Second-half stoppage (0.082) is
the highest rate anywhere --- \textasciitilde1.9× the open-play average.
This is why omitted stoppage minutes cannot be priced at the
match-average rate; but they also cannot simply inherit the stoppage
peak, which is measured over short, unusually urgent windows.}

The three rates that anchor the model (pooled, either-team, per
live-minute):

\begin{longtable}[]{@{}
  >{\raggedright\arraybackslash}p{(\columnwidth - 8\tabcolsep) * \real{0.2000}}
  >{\raggedright\arraybackslash}p{(\columnwidth - 8\tabcolsep) * \real{0.2000}}
  >{\raggedright\arraybackslash}p{(\columnwidth - 8\tabcolsep) * \real{0.2000}}
  >{\raggedright\arraybackslash}p{(\columnwidth - 8\tabcolsep) * \real{0.2000}}
  >{\raggedright\arraybackslash}p{(\columnwidth - 8\tabcolsep) * \real{0.2000}}@{}}
\toprule\noalign{}
\begin{minipage}[b]{\linewidth}\raggedright
Window
\end{minipage} & \begin{minipage}[b]{\linewidth}\raggedright
λ (goals/live-min)
\end{minipage} & \begin{minipage}[b]{\linewidth}\raggedright
95\% Poisson CI
\end{minipage} & \begin{minipage}[b]{\linewidth}\raggedright
goals (n)
\end{minipage} & \begin{minipage}[b]{\linewidth}\raggedright
live-min
\end{minipage} \\
\midrule\noalign{}
\endhead
\bottomrule\noalign{}
\endlastfoot
2nd-half stoppage (decay \textbf{start}) & \textbf{0.0816} & {[}0.0640,
0.1026{]} & 73 & 894.5 \\
1st-half stoppage (used as-is, no decay) & 0.0478 & {[}0.0303, 0.0717{]}
& 23 & 481.2 \\
Regulation open play (decay \textbf{floor}) & \textbf{0.0427} &
{[}0.0396, 0.0461{]} & 675 & 15794.1 \\
\end{longtable}

The 2nd-half stoppage rate is \textbf{1.91×} the open-play floor --- but
it is measured over short, late, high-urgency windows, so we do not
assume it for \emph{all} newly-added minutes. For each marginal omitted
2nd-half minute \texttt{t}, the rate decays geometrically from the
observed rate toward the floor:

\begin{verbatim}
λ(t)    = floor + (obs − floor) · 0.5^(t / h)
λ̄(T,h)  = floor + (obs − floor) · (1 − e^(−kT)) / (kT),   k = ln 2 / h
\end{verbatim}

where \texttt{λ̄(T,h)} is the closed-form \textbf{average} over a window
of \texttt{T} omitted 2nd-half minutes and \texttt{h} is the half-life.
The two natural bounding cases fall out exactly: \texttt{h\ →\ ∞} is
``no decay'' (every omitted minute at the observed rate) and
\texttt{h\ →\ 0} is ``full decay'' (every omitted minute at the floor).
The \textbf{first half keeps its observed rate} --- decay is applied to
the second half only, where the urgency premium is concentrated. Figure
3 shows the decay curves and the rate an average match actually
receives.

\emph{Figure 3. Left: the per-minute rate λ(t) for an omitted 2nd-half
minute, decaying from the observed 2nd-half-stoppage rate (0.0816,
dashed) toward the open-play floor (0.0427, dotted) for half-lives h =
2, 4 (central), 8 minutes; the grey histogram is the distribution of
omitted-minute counts across matches. Right: the window-average rate λ̄ a
match actually receives as a function of its total omitted 2nd-half
minutes T; the green line marks the mean omitted window (T ≈ 6.1 min).
Because most matches have modest omitted windows, the effective rate
sits well above the floor at the central half-life.}

Reported decay band (\texttt{h\ ∈\ {[}2,\ 8{]}} min, central
\texttt{h\ =\ 4}), with the two limits shown for reference:

\begin{longtable}[]{@{}
  >{\raggedright\arraybackslash}p{(\columnwidth - 10\tabcolsep) * \real{0.1667}}
  >{\raggedright\arraybackslash}p{(\columnwidth - 10\tabcolsep) * \real{0.1667}}
  >{\raggedright\arraybackslash}p{(\columnwidth - 10\tabcolsep) * \real{0.1667}}
  >{\raggedright\arraybackslash}p{(\columnwidth - 10\tabcolsep) * \real{0.1667}}
  >{\raggedright\arraybackslash}p{(\columnwidth - 10\tabcolsep) * \real{0.1667}}
  >{\raggedright\arraybackslash}p{(\columnwidth - 10\tabcolsep) * \real{0.1667}}@{}}
\toprule\noalign{}
\begin{minipage}[b]{\linewidth}\raggedright
Full match
\end{minipage} & \begin{minipage}[b]{\linewidth}\raggedright
h=2
\end{minipage} & \begin{minipage}[b]{\linewidth}\raggedright
\textbf{h=4 (central)}
\end{minipage} & \begin{minipage}[b]{\linewidth}\raggedright
h=8
\end{minipage} & \begin{minipage}[b]{\linewidth}\raggedright
h→∞ (no decay)
\end{minipage} & \begin{minipage}[b]{\linewidth}\raggedright
h→0 (full decay)
\end{minipage} \\
\midrule\noalign{}
\endhead
\bottomrule\noalign{}
\endlastfoot
X\% & 22.2\% & \textbf{23.6\%} & 24.9\% & 23.8\% & 17.0\% \\
\end{longtable}

(\texttt{h\ →\ 0} gives 17.0\% on the full match --- not the 9.7\% of
the second-half-only variant --- because the decay floors only the
2nd-half window; the first-half contribution is untouched.)

\textbf{Where it breaks.} If late-game urgency really does persist into
the newly-added minutes, the model \emph{under}-counts and the truth is
nearer 23.8--24.9\%. If the omitted minutes are \emph{deader} than open
play --- plausible, since they are precisely the minutes a referee
declined to play --- even the floor overstates them; but the live-share
term \texttt{ℓ} already strips out dead time, and the gross-up (§4.3)
pushes in the opposite direction. The decay half-life is the explicit
knob that brackets this assumption.

\subsubsection{4.2 The owed-stoppage estimator (formal
definition)}\label{the-owed-stoppage-estimator-formal-definition}

\textbf{Takeaway.} This is the formal construction of \texttt{T\_true},
whose data-quality rationale and calibration to Nate already appear in
§2. It is built bottom-up from dead time we can identify event-by-event,
plus the marker-gated Data Gap allowance and residual constant of §2:

\begin{verbatim}
identifiable = (celebration ∪ substitution ∪ card ∪ injury ∪ excess-restart) ∩ dead-time
T_true       = identifiable + marker-gated Data Gaps + a small residual constant
\end{verbatim}

\texttt{excess-restart} credits only the \emph{excess} of a routine
restart's dead time over a standard allowance (throw-in 20s, goal kick
30s, corner 45s, free kick 60s), so ordinary restarts are not
double-paid. \texttt{marker-gated\ Data\ Gaps} credits a Data Gap (§2)
of ≥20s \textbf{only when its leading edge carries an out-of-play
marker}; the \textasciitilde8.4 min/match of \emph{unmarked} Data Gaps
are dropped and a flat residual constant restores the mean. The
term-by-term build and its calibration to Nate's ``expected'' column (r
= 0.825, MAE 2.44 min) are in §2.

\textbf{Where it breaks.} Only about a quarter of Data Gaps carry the
out-of-play marker, and the source never records ``addability''
directly, so r ≈ 0.82 is effectively the ceiling achievable from free
data. More structurally, the restart allowances and the residual
constant are frozen on \textbf{World Cup 2018} and applied unchanged
where owed time is \textasciitilde1.5--2× as large --- the model's most
exposed assumption (§2, §7).

\subsubsection{4.3 The gross-up: stoppage within
stoppage}\label{the-gross-up-stoppage-within-stoppage}

\textbf{Takeaway.} If you add stoppage minutes, some of \emph{those}
minutes will themselves be eaten by new stoppages, which the directive's
own logic says should also be compensated. The gross-up adds that
second-order time --- but only the genuinely-dead-and-addable fraction
\texttt{z} of added time recurs, not the whole non-live share, so the
effect is small and is \textbf{not} double-counting.

\texttt{z} is measured directly: of 44.2 dead minutes per match in
regulation, the estimator counts 16.9 as genuine addable stoppage, so
\texttt{z\ =\ 16.9\ /\ 44.2\ =\ 0.382}. The one-pass live factor is
\texttt{φ(s)\ =\ s·(1\ +\ z(1−s))}; iterating stoppage-within-stoppage
to a fixed point gives the geometric \texttt{s\ /\ (1\ −\ z(1−s))}.
Because the recurring ratio is \texttt{z(1−s)\ ≈\ 0.18} --- not the
naive non-live share \texttt{1−s\ ≈\ 0.46} --- the geometric tail
collapses to just above the single pass:

\begin{longtable}[]{@{}llll@{}}
\toprule\noalign{}
Full match & gross-up OFF & \textbf{ON (central)} & geometric ceiling \\
\midrule\noalign{}
\endhead
\bottomrule\noalign{}
\endlastfoot
X\% & 21.1\% & \textbf{23.6\%} & 24.2\% \\
\end{longtable}

\textbf{Where it breaks.} If you reject second-order compensation
entirely, X\% falls only to \textbf{21.1\%} (still about 1 in 5). If you
believed \emph{all} dead time recurs (\texttt{z\ =\ 1}), the ceiling
would be \texttt{1/s\ ≈\ 34\%} --- the \texttt{z}-weighting is exactly
what kills that tail. Either way, the whole gross-up question moves the
headline by ≤ 3 points.

(The model's fourth load-bearing judgment --- how the ambiguous Data
Gaps are resolved, and the calibration that pins them to Nate --- is in
§2, the section flagged for review. Its propagated uncertainty, the
estimator's calibration error, enters the bootstrap as a per-match term
σ = MAE·√(π/2) ≈ 3.06 min under a half-normal; see §5.)

\begin{center}\rule{0.5\linewidth}{0.5pt}\end{center}

\subsection{5. Uncertainty --- sampling vs
specification}\label{uncertainty-sampling-vs-specification}

\textbf{Takeaway.} Two distinct kinds of uncertainty are kept separate.
\textbf{(a) Sampling}: a bootstrap CI reflecting finite goal counts and
estimator error. \textbf{(b) Specification}: how X\% moves as modeling
choices change. Once the Data Gap handling is fixed (§2), the
specification spread is comparable to the sampling spread --- neither
dominates.

\textbf{(a) Sampling CI = {[}20.6\%, 27.4\%{]}} (width 6.7 pts). In each
of B = 1000 bootstrap draws (seed fixed): every distinct λ-cell is
redrawn from its Jeffreys--Gamma posterior
\texttt{Gamma(count\ +\ 0.5,\ 1/exposure)} --- the 2nd-half decay
redraws \emph{both} the 73-goal observed cell and the 675-goal floor
cell and recombines them through λ̄ --- plus one shared owed-stoppage
estimator-error draw \texttt{N(0,\ 3.06\ min)}, split across the two
halves by their share of addable time and re-shifting the decay horizon.
This is genuine sampling uncertainty in the inputs.

\textbf{(b) Specification sensitivities} (each choice swept with the
others held at central, full match):

\begin{longtable}[]{@{}
  >{\raggedright\arraybackslash}p{(\columnwidth - 4\tabcolsep) * \real{0.3333}}
  >{\raggedright\arraybackslash}p{(\columnwidth - 4\tabcolsep) * \real{0.3333}}
  >{\raggedright\arraybackslash}p{(\columnwidth - 4\tabcolsep) * \real{0.3333}}@{}}
\toprule\noalign{}
\begin{minipage}[b]{\linewidth}\raggedright
Modeling choice
\end{minipage} & \begin{minipage}[b]{\linewidth}\raggedright
Levels → X\%
\end{minipage} & \begin{minipage}[b]{\linewidth}\raggedright
Spread
\end{minipage} \\
\midrule\noalign{}
\endhead
\bottomrule\noalign{}
\endlastfoot
\textbf{λ source} & all-pooled \textbf{23.6} · POST-only 22.6 ·
regime-matched 23.8 · PRE-only 26.1 & \textasciitilde3.5 pts \\
\textbf{Gross-up} & off 21.1 → \textbf{on 23.6} → geometric 24.2 &
\textasciitilde3.1 pts \\
\textbf{Decay half-life} & h=2 22.2 · \textbf{h=4 23.6} · h=8 24.9 &
\textasciitilde2.7 pts \\
\textbf{Conditioning} & overall \textbf{23.6} · split by tied/not-tied
23.4 & \textasciitilde0.2 pts (negligible) \\
One-at-a-time band (min--max of the above) & \textbf{21.1\% -- 26.1\%} &
5.0 pts ≈ 0.7× sampling \\
Full \textbf{joint} envelope (all choices together) & \textbf{18.6\% --
27.3\%} & 8.7 pts ≈ 1.3× sampling \\
\emph{Data Gap handling, if (improperly) swept as a knob} & \emph{10.8\%
-- 37.3\%} & \emph{\textasciitilde26 pts --- see §2} \\
\end{longtable}

The \textbf{λ source} row is where the PRE/POST split (§2) earns its
keep. The 2022 directive raised stoppage \emph{played} sharply --- 6.8 →
10.0 min/match, PRE → POST --- yet goals \emph{per live-minute} barely
moved: pricing omitted minutes at PRE-only, POST-only, or pooled rates
spans just 22.6--26.1\%, all within Poisson noise. Stoppage time
changed; productivity did not. That is precisely why the model pools all
six tournaments for λ (``all-pooled'' above), and why the top-end 26.1\%
comes entirely from the thin, wide-CI PRE-only cohort rather than any
real PRE/POST productivity gap. (``regime-matched'' prices PRE omitted
minutes at the PRE rate and POST minutes at the POST rate, rather than
pooling.)

\textbf{Which choices move X\% most.} The \textbf{λ source} (driven
entirely by the thin, wide-CI PRE-only cohort at the top end) and the
\textbf{gross-up} are the widest legitimate axes, \textasciitilde3
points each; the \textbf{decay half-life} is \textasciitilde2.7;
conditioning is noise. Both extremes of the reported envelope are
identifiable: the floor (18.6--21.1\%) is the gross-up-off / most-decay
corner; the ceiling (26.1--27.3\%) is the PRE-only λ × least-decay
corner --- i.e.~the least-certain input cohort. (The envelope sweeps the
half-life over its plausible 2--8 min range; the two mathematical limits
--- no decay, full decay --- are shown in §4.1 only to recover the
bounding rails and are not part of the band.)

\textbf{A structural soft spot.} The first-half window contributes
\textbf{+7.6 points} (second-half-only 16.0\% → full-match 23.6\%) ---
roughly a third of the headline --- and rests on the \textbf{23-goal}
first-half-stoppage rate. The bootstrap propagates that thinness (a wide
first-half CI), but the \emph{point} estimate's dependence on a 23-goal
cell is worth naming.

\begin{center}\rule{0.5\linewidth}{0.5pt}\end{center}

\subsection{6. External validation}\label{external-validation}

\textbf{Takeaway.} Two external anchors certify the model's
\textbf{inputs}; neither certifies its \textbf{output}.

\begin{itemize}
\tightlist
\item
  \textbf{Nate Silver, World Cup 2018 (Figure 1, §2).} Detailed in §2:
  our owed-stoppage estimator tracks his hand-measured ``expected'' at
  \textbf{r = 0.825} (MAE 2.44 min) and our stoppage-played clock tracks
  his ``actual'' at \textbf{r = 0.992}. The point worth repeating here
  is that the owed-vs-played \emph{gap} is not a modeling artifact ---
  Nate, with a stopwatch and no access to our pipeline, independently
  found \textbf{13.16 owed vs 6.98 played}, i.e.~roughly half the owed
  stoppage going unplayed.
\item
  \textbf{Opta, World Cup 2022 ball-in-play.} Opta publishes an official
  average ball-in-play time of 58:04 per match for World Cup 2022; our
  gap-method reconstruction gives 57:40 (24s low). This certifies the
  live-share denominator \texttt{s}.
\end{itemize}

\textbf{What they do \emph{not} certify.} Both anchors are World Cup
2018/2022 respectively; the POST owed-time magnitude
(\textasciitilde19.4 min) is an extrapolation.

\subsection{7. Limitations \& open
questions}\label{limitations-open-questions}

Honest list, least-defensible first:

\begin{enumerate}
\def\labelenumi{\arabic{enumi}.}
\tightlist
\item
  \textbf{No independent counterfactual benchmark.} Every external check
  is on an input; X\% itself is unfalsified against any outside number.
\item
  \textbf{POST coverage gap.} The estimator is calibrated on World Cup
  2018 only. The POST cohort --- where the headline mostly lives, and
  where owed time is \textasciitilde1.5--2× World Cup 2018 --- rides on
  frozen-2018 constants and a single World Cup 2022 ball-in-play point.
  If the 2018 marker/allowance/residual structure does not transfer to
  Copa América and AFCON, the POST \texttt{T\_true} (and thus X\%) is
  the most exposed quantity.
\item
  \textbf{Data-Gap-as-point (§2).} Defensible, but it is the choice that
  makes the band tight; a reviewer who insists ``none''/``all'' are
  legitimate uncertainty will see a far wider honest range.
\item
  \textbf{First-half thinness (§5).} About a third of the headline rests
  on 23 goals, with no game-state propagation across the first-half
  increment (A1).
\item
  \textbf{Productivity transfer is an assumption, not a measurement
  (§4.1).} The decay band brackets the urgency premium but cannot rule
  out that \emph{omitted} minutes --- the ones a referee chose not to
  play --- are systematically unlike any observed minute.
\item
  \textbf{Cooling/water breaks} were investigated as an additional
  stoppage source and set aside: roughly three-quarters of that time is
  already captured by the estimator, and crediting the rest added no
  robust accuracy on the only externally-validated set.
\end{enumerate}

\end{document}
